\section{Introduction}\label{sec:intro}

Identified dynamical models are not built once and discarded: they are simulated over long horizons, embedded in
predictive controllers and estimators, and executed on processors with limited energy budgets. Each evaluation
consumes arithmetic operations, and the energy and emissions associated with them accumulate over the life of the
model \cite{PWG2018,LGI2021}. For polynomial NARX (nonlinear autoregressive with exogenous inputs) models, this
per-sample cost is fixed by the structure: every retained regressor adds an addition, and every factor adds a
multiplication. Green-Box system identification \cite{NMN2023,NZBN2026} makes this cost explicit through an
arithmetic-complexity measure that weights additions and multiplications by a bit-level cost model, and treats it as
a design objective alongside accuracy. Selecting the structure of such a model, that is, which delayed inputs,
outputs and products enter it, is therefore a problem with two competing objectives rather than a pure accuracy
problem.

Structure selection for NARX models is classically performed by forward regression with orthogonal least squares,
which ranks candidate regressors by their error reduction ratio \cite{KBLM1988,Bil2013}, and the Green-Box cost has
been inserted into this greedy procedure both when ranking regressors and when choosing the model size
\cite{NMN2026GB}. Evolutionary methods explore structures combinatorially instead. Multi-objective genetic
programming (GP) has identified NARX structures with model size as a second objective \cite{RFF2004}, GP combined with
orthogonal least squares has identified input--output models \cite{MAS2005}, and multi-objective evolutionary
frameworks for nonlinear identification have been compared systematically \cite{HSM2020}. Multigene GP (MGGP)
\cite{HHM+1996,SLW2010} is well suited to polynomial NARX models: with multiplication as the only arithmetic
primitive and delays introduced by backshift operators, each gene is a monomial and each individual is a model whose
coefficients follow from least squares. The open-source \texttt{mggp} package \cite{mggpy} implements this
representation and has recently added a bi-objective mode in which the user supplies the second objective.

Three requirements remain open between these lines of work. In Pareto GP, complexity is usually measured by the number
of nodes or terms, or by the order of non-linearity \cite{SK2005,VSH2009}, and sparse-regression methods measure it
by the number of active terms \cite{BPK2016}; none of these measures separates the cost of an addition from that of a
multiplication, nor accounts for regressors that are duplicated in the representation but executed once. The
Green-Box measure does separate them, but it fixes their relative weight by a bit-level model of multiplication, whereas
the cost actually incurred depends on the platform that executes the model, and this weight has not been compared
with measured execution. Moreover, supplying a second objective to an evolutionary algorithm does not by itself make
the search bi-objective: when selection compares fitness tuples lexicographically, the second objective only breaks
ties of the first, and the trade-off is recovered only from whatever the final population contains. The question
addressed in this paper is therefore whether an operation-level cost, computed on the executed polynomial and weighted
for the executing platform, can drive MGGP towards better accuracy--cost trade-offs, how much of this depends on the
selection scheme rather than on the objective itself, and whether the resulting cost reductions carry over to the
measured time and energy of executing the identified models.

This paper proposes Green MGGP, which uses the arithmetic cost of the canonical, duplicate-free polynomial as the
second objective of MGGP and replaces lexicographic selection by Pareto-based selection with crowding and structural
duplicate control, while keeping the representation, operators and estimators of the \texttt{mggp} package. The cost
counts additions and multiplications with a multiplication weight $\rho$ that represents the executing platform; it
contains the bit-level Green-Box complexity as the special case $\rho_K=B^{\alpha-1}$ and the operation count as $\rho=1$.
A rule that chooses one model from the resulting set using identification data only is also defined. The formulation
is supported by four results: the identified model, and hence every objective, depends on an individual only through
its set of distinct monomials; the arithmetic cost is a strictly monotone, modular function of the numbers of
additions and multiplications for every weight; the lexicographic mode reproduces the single-objective search exactly
unless errors tie between individuals of different cost; and, under a capacity condition, the Pareto-based scheme
never lets the population front deteriorate. The method is compared with the single-objective package and with its
lexicographic bi-objective mode on the regression examples distributed with the package, namely a bilinear system with
and without measurement noise and a Bouc--Wen hysteretic system, over \nSeeds{} independent runs and under both
weights, and the free-run simulation of the selected models is timed and its energy measured. The measured execution
times are explained by the number of operations, with multiplications costing about as much as additions, which
supports $\rho=1$ on the platform used. With this weight, Green MGGP yields higher identification and free-run
validation hypervolumes than both baselines in every run on all three examples, its selected models need
$17$--$32\%$ fewer operations than the single-objective ones, and its trade-off set always contains
the true regressors of the noisy system; the same holds for the hypervolume under the bit-level weight. Executed as
canonical polynomials, the selected Green models required $13$--$35\%$ less energy than the single-objective ones.
These results show that the selection scheme, not only the cost function, determines whether the arithmetic cost shapes
the identified structures, and that the weight of the cost should be calibrated on the target platform.
